\begin{derivation}[从\eqnrefs{eq:born_nonmarkov_redfield_kernel}到\eqnrefs{eq:born_nonmarkov_noise_dissipation_form}]
记

\begin{equation*}
A=q(t),\qquad B_\tau=q(t-\tau),\qquad \rho_\tau=\rho_{S,I}(t-\tau),
\end{equation*}
并把平稳环境相关函数写成
\begin{equation*}
C(\tau)=C_{\rm sym}(\tau)+\ii C_{\rm asym}(\tau),\qquad
C(-\tau)=C_{\rm sym}(\tau)-\ii C_{\rm asym}(\tau).
\end{equation*}
二阶 Born 核中的算符组合为
\begin{align*}
\mathcal K_\tau
={}&C(\tau)[A,B_\tau\rho_\tau]
+C(-\tau)[\rho_\tau B_\tau,A]\\
={}&(C_{\rm sym}+\ii C_{\rm asym})(AB_\tau\rho_\tau-B_\tau\rho_\tau A)
+(C_{\rm sym}-\ii C_{\rm asym})(\rho_\tau B_\tau A-A\rho_\tau B_\tau).
\end{align*}
先收集 $C_{\rm sym}$ 项：
\begin{align*}
\mathcal K_\tau^{({\rm sym})}
&=C_{\rm sym}(AB_\tau\rho_\tau-A\rho_\tau B_\tau-B_\tau\rho_\tau A+\rho_\tau B_\tau A)\\
&=C_{\rm sym}[A,[B_\tau,\rho_\tau]].
\end{align*}
再收集 $C_{\rm asym}$ 项：
\begin{align*}
\mathcal K_\tau^{({\rm asym})}
&=\ii C_{\rm asym}(AB_\tau\rho_\tau+A\rho_\tau B_\tau-B_\tau\rho_\tau A-\rho_\tau B_\tau A)\\
&=\ii C_{\rm asym}[A,\{B_\tau,\rho_\tau\}].
\end{align*}
因此
\begin{equation*}
\mathcal K_\tau
=C_{\rm sym}(\tau)[q(t),[q(t-\tau),\rho_\tau]]
+\ii C_{\rm asym}(\tau)[q(t),\{q(t-\tau),\rho_\tau\}].
\end{equation*}
把它代回二阶 Born 方程前面的 $-\lambda^2/\hbar^2$ 系数和 $\tau$ 积分，就得到\eqnrefs{eq:born_nonmarkov_noise_dissipation_form}。
\end{derivation}

\begin{derivation}[从\eqnrefs{eq:generalized_langevin_from_influence}到\eqnrefs{eq:dressed_oscillator_retarded_response}]
对谐振子势 $V(q)=M\Omega_0^2q^2/2$，正文\eqnrefs{eq:generalized_langevin_from_influence} 可写成
\[
M\ddot q(t)+M\Omega_0^2q(t)
+\int_{-\infty}^{t}\dd s\,\chi_B^R(t-s)q(s)
=\xi(t).
\]
在推迟核中令 $\tau=t-s$，则
\[
\int_{-\infty}^{t}\dd s\,\chi_B^R(t-s)q(s)
=\int_0^\infty\dd\tau\,\chi_B^R(\tau)q(t-\tau).
\]
采用正文的 Fourier 约定
$q(t)=\int \dd\omega\,q(\omega)\ee^{-\ii\omega t}/(2\pi)$，卷积项变成
\begin{align*}
&\int_0^\infty\dd\tau\,\chi_B^R(\tau)
\int\frac{\dd\omega}{2\pi}q(\omega)
\ee^{-\ii\omega(t-\tau)}\\
&=\int\frac{\dd\omega}{2\pi}
\left[
\int_0^\infty\dd\tau\,
\chi_B^R(\tau)\ee^{+\ii\omega\tau}
\right]q(\omega)\ee^{-\ii\omega t}\\
&=\int\frac{\dd\omega}{2\pi}
\chi_B^R(\omega)q(\omega)\ee^{-\ii\omega t}.
\end{align*}
同时 $\ddot q(t)\mapsto-\omega^2q(\omega)$。因此每个频率分量满足
\[
\left[
M\bigl(\Omega_0^2-(\omega+\ii0^+)^2\bigr)
+\chi_B^R(\omega)
\right]q(\omega)=\xi(\omega).
\]
对乘在 $q(\omega)$ 前的标量核求逆，得到
\[
q(\omega)=\mathcal R_q(\omega)\xi(\omega),
\qquad
\mathcal R_q(\omega)
=\frac{1}{M[\Omega_0^2-(\omega+\ii0^+)^2]+\chi_B^R(\omega)},
\]
即正文\eqnrefs{eq:dressed_oscillator_retarded_response}。
\end{derivation}

\begin{derivation}[从\eqnrefs{eq:dressed_coordinate_spectrum_dp68}到\eqnrefs{eq:dressed_system_fdt_from_bath_fdt}]

正文\eqnrefs{eq:dressed_coordinate_spectrum_dp68} 表明系统坐标谱等于响应模方乘环境噪声谱。热平衡时，环境噪声谱又由\eqnrefs{eq:fdt_from_defined_correlators} 与环境响应的虚部关联。利用复数逆函数的虚部恒等式，把 $|\mathcal R_{x_S}|^2[-\operatorname{Im}\chi_B^R]$ 改写成系统自身响应的耗散部分，就得到\eqnrefs{eq:dressed_system_fdt_from_bath_fdt}。

令
\[
D(\omega)=M[\Omega_0^2-(\omega+\ii0^+)^2]+\chi_B^R(\omega),
\qquad
\mathcal R_{x_S}=D^{-1}.
\]
于是
\[
\operatorname{Im}\mathcal R_{x_S}
=-|\mathcal R_{x_S}|^2\operatorname{Im}\chi_B^R.
\]
按正文约定 $\chi_{x_S}^R=-\mathcal R_{x_S}$，因此
\[
-\operatorname{Im}\chi_{x_S}^R
=|\mathcal R_{x_S}|^2[-\operatorname{Im}\chi_B^R].
\]
将这一恒等式和环境 FDT~\eqnrefs{eq:fdt_from_defined_correlators} 代入\eqnrefs{eq:dressed_coordinate_spectrum_dp68}，得到\eqnrefs{eq:dressed_system_fdt_from_bath_fdt}。
\end{derivation}

% 数学结论及完整证明已移入正文相应公式之后。


% 数学结论及完整证明已移入正文相应公式之后。


% 数学结论及完整证明已移入正文相应公式之后。


% 数学结论及完整证明已移入正文相应公式之后。


% 数学结论及完整证明已移入正文相应公式之后。


% 数学结论及完整证明已移入正文相应公式之后。


% 数学结论及完整证明已移入正文相应公式之后。


% 数学结论及完整证明已移入正文相应公式之后。


% 数学结论及完整证明已移入正文相应公式之后。


% 数学结论及完整证明已移入正文相应公式之后。


% 数学结论及完整证明已移入正文相应公式之后。


\begin{derivation}[从\eqnrefs{eq:exact_second_order_integral_identity_open_system}到\eqnrefs{eq:born_nonmarkov_redfield_kernel}]
相互作用绘景中的联合态满足
\[
\dot\rho_I(t)=-\frac{\ii}{\hbar}[H_I(t),\rho_I(t)],
\qquad
H_I(t)=\lambda\sum_\alpha A_\alpha(t)\otimes B_\alpha(t).
\]
从 $0$ 积分到 $t$：
\[
\rho_I(t)=\rho_I(0)
-\frac{\ii}{\hbar}\int_0^t\dd s\,[H_I(s),\rho_I(s)].
\]
把这个积分式代回原微分方程，并令 $s=t-\tau$，得到正文\eqnrefs{eq:exact_second_order_integral_identity_open_system}：
\[
\dot\rho_I(t)
=-\frac{\ii}{\hbar}[H_I(t),\rho_I(0)]
-\frac1{\hbar^2}\int_0^t\dd\tau\,
[H_I(t),[H_I(t-\tau),\rho_I(t-\tau)]].
\]
取初态因子化 $\rho_I(0)=\rho_S(0)\otimes\rho_B$，并令
$\Tr_B(B_\alpha\rho_B)=0$。对环境取迹以后，一阶项为
\begin{align*}
\Tr_B[H_I(t),\rho_S(0)\otimes\rho_B]
&=\lambda\sum_\alpha
[A_\alpha(t),\rho_S(0)]\Tr_B[B_\alpha(t)\rho_B]\\
&=0.
\end{align*}
Born 截断只用于二阶积分核：
\[
\rho_I(t-\tau)
=\rho_{S,I}(t-\tau)\otimes\rho_B+O(\lambda).
\]
由于核前已有两个 $H_I\propto\lambda$，代入后被丢弃部分从 $O(\lambda^3)$ 开始。

把双重对易子完整展开，记
$\rho_S^-=\rho_{S,I}(t-\tau)$：
\begin{align*}
&[H_I(t),[H_I(t-\tau),\rho_S^-\otimes\rho_B]]\\
={}&H_I(t)H_I(t-\tau)(\rho_S^-\otimes\rho_B)
-H_I(t)(\rho_S^-\otimes\rho_B)H_I(t-\tau)\\
&-H_I(t-\tau)(\rho_S^-\otimes\rho_B)H_I(t)
+(\rho_S^-\otimes\rho_B)H_I(t-\tau)H_I(t).
\end{align*}
代入 $H_I=\lambda\sum A_\alpha\otimes B_\alpha$。第一项的环境偏迹为
\begin{align*}
&\lambda^2\sum_{\alpha\beta}
A_\alpha(t)A_\beta(t-\tau)\rho_S^-\,
\Tr_B[B_\alpha(t)B_\beta(t-\tau)\rho_B]\\
&=\lambda^2\sum_{\alpha\beta}
C_{\alpha\beta}(\tau)
A_\alpha(t)A_\beta(t-\tau)\rho_S^-.
\end{align*}
平稳性使相关函数只依赖时间差。第二项利用迹的循环性：
\begin{align*}
&\Tr_B[B_\alpha(t)\rho_BB_\beta(t-\tau)]\\
&=\Tr_B[B_\beta(t-\tau)B_\alpha(t)\rho_B]
=C_{\beta\alpha}(-\tau),
\end{align*}
故给
\[
-\lambda^2\sum_{\alpha\beta}
C_{\beta\alpha}(-\tau)
A_\alpha(t)\rho_S^-A_\beta(t-\tau).
\]
第三、第四项同理分别为
\begin{align*}
&-\lambda^2\sum_{\alpha\beta}
C_{\alpha\beta}(\tau)
A_\beta(t-\tau)\rho_S^-A_\alpha(t),\\
&+\lambda^2\sum_{\alpha\beta}
C_{\beta\alpha}(-\tau)
\rho_S^-A_\beta(t-\tau)A_\alpha(t).
\end{align*}
把第一项与第三项组合成一个对易子，把第二项与第四项组合成另一个对易子：
\begin{align*}
\Tr_B[H_I(t),[H_I(t-\tau),\rho_S^-\otimes\rho_B]]
=\lambda^2\sum_{\alpha\beta}\Big\{&
C_{\alpha\beta}(\tau)
[A_\alpha(t),A_\beta(t-\tau)\rho_S^-]\\
+&C_{\beta\alpha}(-\tau)
[\rho_S^-A_\beta(t-\tau),A_\alpha(t)]\Big\}.
\end{align*}
代回二次积分恒等式后得到
\begin{align*}
\dot\rho_{S,I}(t)
={}&-\frac{\lambda^2}{\hbar^2}
\sum_{\alpha\beta}\int_0^t\dd\tau\,\Big\{
C_{\alpha\beta}(\tau)
[A_\alpha(t),A_\beta(t-\tau)\rho_{S,I}(t-\tau)]\\
&\hspace{30mm}+C_{\beta\alpha}(-\tau)
[\rho_{S,I}(t-\tau)A_\beta(t-\tau),A_\alpha(t)]\Big\}
+O(\lambda^3),
\end{align*}
即正文\eqnrefs{eq:born_nonmarkov_redfield_kernel}。
\end{derivation}

% 数学结论及完整证明已移入正文相应公式之后。


% 数学结论及完整证明已移入正文相应公式之后。


% 数学结论及完整证明已移入正文相应公式之后。


% 数学结论及完整证明已移入正文相应公式之后。


% 数学结论及完整证明已移入正文相应公式之后。


% 数学结论及完整证明已移入正文相应公式之后。


\begin{derivation}[从\eqnrefs{eq:gaussian_feynman_vernon_influence}到\eqnrefs{eq:generalized_langevin_from_influence}]
取线性耦合 $A=x_S$。要从影响泛函得到运动方程，必须先把它的两个核分成“噪声权重”和“因果反作用”两部分，而不能直接把指数读成 Langevin 方程。

首先利用 $C_{\rm sym}(t,s)=C_{\rm sym}(s,t)$。正文影响泛函中的对称核积分只覆盖 $t>s$ 的三角域，因此
\begin{align*}
 &\int_{t_i}^{t_f}\dd t\int_{t_i}^{t}\dd s\,
 x_\Delta(t)C_{\rm sym}(t,s)x_\Delta(s)\\
 &\qquad=\frac12
 \int_{t_i}^{t_f}\dd t\int_{t_i}^{t_f}\dd s\,
 x_\Delta(t)C_{\rm sym}(t,s)x_\Delta(s).
\end{align*}
于是这一部分恰好是高斯特征泛函
\begin{align*}
 &\exp\!\left[-\frac{1}{2\hbar^2}
 \iint\dd t\dd s\,
 x_\Delta(t)C_{\rm sym}(t,s)x_\Delta(s)\right]\\
 &\qquad=
 \int\mathcal D\xi\,\mathcal P[\xi]
 \exp\!\left[\frac{\ii}{\hbar}
 \int\dd t\,\xi(t)x_\Delta(t)\right],
\end{align*}
其中 $\mathcal P[\xi]$ 是零均值高斯测度，其协方差由
\begin{equation*}
 \langle\xi(t)\xi(s)\rangle_\xi=C_{\rm sym}(t,s)
\end{equation*}
唯一固定。这里的 $\xi$ 首先只是 Hubbard--Stratonovich 辅助场；只有在相应半经典极限中，才把它解释成一条经典随机力轨迹。

再处理反对称核。由正文定义
\begin{equation*}
 \chi_B^R(t,s)
 =\frac{2}{\hbar}\Theta(t-s)C_{\rm asym}(t,s),
\end{equation*}
所以影响泛函中的相位项精确改写为
\begin{align*}
 &-\frac{2\ii}{\hbar^2}
 \int_{t>s}\dd t\dd s\,
 x_\Delta(t)C_{\rm asym}(t,s)x_\Sigma(s)\\
 &\qquad=-\frac{\ii}{\hbar}
 \int_{t_i}^{t_f}\dd t\,x_\Delta(t)
 \int_{t_i}^{t}\dd s\,
 \chi_B^R(t,s)x_\Sigma(s).
\end{align*}
阶跃函数使积分只依赖 $s\le t$ 的过去历史，这就是“记忆力是因果的”的数学来源。

系统自身的作用量取
\begin{equation*}
 S_S[x]=\int\dd t\left[\frac{M}{2}\dot x^2-V(x)\right],
\end{equation*}
并写 $x_\pm=x_\Sigma\pm x_\Delta/2$，则动能差为
\begin{align*}
 \frac{M}{2}(\dot x_+^2-\dot x_-^2)
 &=M\dot x_\Sigma\dot x_\Delta\\
 &=-M\ddot x_\Sigma x_\Delta
 +\frac{\dd}{\dd t}(M\dot x_\Sigma x_\Delta),
\end{align*}
而势能差的 Taylor 展开给
\begin{equation*}
 V(x_+)-V(x_-)
 =V'(x_\Sigma)x_\Delta+O(x_\Delta^3).
\end{equation*}
在两端固定路径差 $x_\Delta(t_i)=x_\Delta(t_f)=0$ 时，边界项消失，因此
\begin{equation*}
 S_S[x_+]-S_S[x_-]
 =\int\dd t\,x_\Delta(t)
 [-M\ddot x_\Sigma(t)-V'(x_\Sigma(t))]
 +O(x_\Delta^3).
\end{equation*}

把系统作用量、耗散相位和辅助高斯场合并，总有效指数中关于 $x_\Delta$ 的线性项为
\begin{align*}
 \frac{\ii}{\hbar}\int\dd t\,x_\Delta(t)
 \Bigg[&-M\ddot x_\Sigma(t)-V'(x_\Sigma(t))\\
 &-\int_{t_i}^{t}\dd s\,\chi_B^R(t,s)x_\Sigma(s)
 +\xi(t)\Bigg].
\end{align*}
粗粒化的半经典驻值方程就是令这个线性系数对每个 $t$ 都为零。最后在物理路径上令 $x_+=x_-=x_S$，即 $x_\Delta=0$、$x_\Sigma=x_S$，得到
\begin{equation*}
 M\ddot x_S(t)+V'[x_S(t)]
 +\int_{t_i}^{t}\dd s\,\chi_B^R(t,s)x_S(s)
 =\xi(t),
\end{equation*}
并同时保留
\begin{equation*}
 \langle\xi(t)\xi(s)\rangle=C_{\rm sym}(t,s).
\end{equation*}
这才是正文\eqnrefs{eq:generalized_langevin_from_influence}。推导清楚地区分了三个层次：$C_{\rm asym}$ 通过推迟核生成确定性的历史反作用；$C_{\rm sym}$ 通过高斯特征泛函生成涨落；把辅助场 $\xi$ 当作经典随机力还需要额外的半经典解释，而不是影响泛函本身自动给出的本体论结论。
\end{derivation}

% 数学结论及完整证明已移入正文相应公式之后。


% 数学结论及完整证明已移入正文相应公式之后。


% 数学结论及完整证明已移入正文相应公式之后。


\begin{derivation}[从\eqnrefs{eq:open_jc_effective_master_equation}到\eqnrefs{eq:open_jc_cooperativity}]
先从主方程读取电子相干线宽。对 $\rho_{ge}=\langle g|\rho|e\rangle$，布居弛豫耗散子给
\begin{align*}
\langle g|\mathcal D[\sigma_-]\rho|e\rangle
&=\langle g|\sigma_-\rho\sigma_+|e\rangle
-\frac12\langle g|\{\sigma_+\sigma_-,\rho\}|e\rangle\\
&=-\frac12\rho_{ge},
\end{align*}
而
\begin{equation*}
\langle g|\mathcal D[\sigma_z]\rho|e\rangle
=\langle g|\sigma_z\rho\sigma_z-\rho|e\rangle
=-2\rho_{ge}.
\end{equation*}
所以
\begin{equation*}
\dot\rho_{ge}\big|_{\rm diss}
=-\left(\frac{\Gamma_1}{2}+\Gamma_\varphi\right)\rho_{ge}
\equiv-\Gamma_2\rho_{ge}.
\end{equation*}

在线性单激发子空间 $\{|e,0\rangle,|g,1\rangle\}$ 中，零失谐时的均值方程为
\begin{equation*}
\frac{\dd}{\dd t}
\begin{pmatrix}c_e\\c_c\end{pmatrix}
=
\underbrace{\begin{pmatrix}
-\Gamma_2&-\ii g_c\\
-\ii g_c^*&-\Gamma_{\rm cav}/2
\end{pmatrix}}_{\mathsf M}
\begin{pmatrix}c_e\\c_c\end{pmatrix}.
\end{equation*}
矩阵 $\mathsf M$ 的两个本征值由二次方程精确给出
\begin{equation*}
\lambda_\pm
=-\frac{\Gamma_2+\Gamma_{\rm cav}/2}{2}
\pm
\sqrt{
\frac{(\Gamma_2-\Gamma_{\rm cav}/2)^2}{4}
-|g_c|^2}.
\end{equation*}
因此时间域出现振荡型正常模分裂的条件是根号变成纯虚数，即
\begin{equation*}
|g_c|>
\frac12\left|\Gamma_2-\frac{\Gamma_{\rm cav}}{2}\right|.
\end{equation*}
这是一条关于本征值结构的判据；实验上要把两条谱峰清楚分辨，还需分裂大于相应总线宽。

另一方面，在坏腔极限中可从第二行写出完全等价的恒等式
\begin{equation*}
c_c
=-\frac{2\ii g_c^*}{\Gamma_{\rm cav}}c_e
-\frac{2}{\Gamma_{\rm cav}}\dot c_c.
\end{equation*}
把它代回第一行，不丢项时得到
\begin{equation*}
\dot c_e
=-\left(\Gamma_2+\frac{2|g_c|^2}{\Gamma_{\rm cav}}\right)c_e
+\frac{2\ii g_c}{\Gamma_{\rm cav}}\dot c_c.
\end{equation*}
当快变量条件使最后一项相对于前两项可忽略时，腔介导的电子振幅衰减率为
\begin{equation*}
\Gamma_{\rm cav\to e}^{({\rm amp})}
\equiv\frac{2|g_c|^2}{\Gamma_{\rm cav}}.
\end{equation*}
将它与电子自身相干衰减率比较，严格得到正文采用的协同性定义
\begin{equation*}
\frac{\Gamma_{\rm cav\to e}^{({\rm amp})}}{\Gamma_2}
=\frac{2|g_c|^2}{\Gamma_{\rm cav}\Gamma_2}
\equiv\mathscr C_{\rm coop}.
\end{equation*}
\end{derivation}

% 数学结论及完整证明已移入正文相应公式之后。


% 数学结论及完整证明已移入正文相应公式之后。


% 数学结论及完整证明已移入正文相应公式之后。


% 数学结论及完整证明已移入正文相应公式之后。


% 数学结论及完整证明已移入正文相应公式之后。


% 数学结论及完整证明已移入正文相应公式之后。


% 数学结论及完整证明已移入正文相应公式之后。


% 数学结论及完整证明已移入正文相应公式之后。


% 数学结论及完整证明已移入正文相应公式之后。


% 数学结论及完整证明已移入正文相应公式之后。


% 数学结论及完整证明已移入正文相应公式之后。


% 数学结论及完整证明已移入正文相应公式之后。


% 数学结论及完整证明已移入正文相应公式之后。


\begin{derivation}[从自由传播核到高斯解，并直接代回 TDSE]
在二阶包络模型中，连续动量谱给传播核
\begin{equation*}
 K(\zeta,\zeta';t)=\int\frac{\dd p}{2\pi\hbar}
 \exp\!\left[\frac{\ii p(\zeta-\zeta')}{\hbar}
             -\frac{\ii p^2t}{2m_\parallel\hbar}\right]
 =\sqrt{\frac{m_\parallel}{2\pi\ii\hbar t}}
 \exp\!\left[\frac{\ii m_\parallel(\zeta-\zeta')^2}{2\hbar t}\right],
 \quad t>0.
\end{equation*}
振荡高斯积分可先加正的收敛因子，再取其趋零极限。将初态代入 $\psi(\zeta,t)=\int K\psi(\zeta',0)\dd\zeta'$，整理为
\begin{equation*}
 -A\zeta'^2+B\zeta'+C,
 \quad A=a_0-\frac{\ii m_\parallel}{2\hbar t},\quad
 B=-\frac{\ii m_\parallel\zeta}{\hbar t},\quad
 C=\frac{\ii m_\parallel\zeta^2}{2\hbar t}.
\end{equation*}
$\Re A>0$ 保证积分收敛；使用 $\int\ee^{-Ax^2+Bx}\dd x=\sqrt{\pi/A}\ee^{B^2/(4A)}$，合并前因子与指数即得\eqnrefs{eq:electron-gaussian-tdse-solution}。

还可独立代回检验。写 $a(t)=a_0/d(t)$、$\mathcal N(t)=\mathcal N_0/\sqrt{d(t)}$，则
\begin{equation*}
 \dot a=-\frac{2\ii\hbar}{m_\parallel}a^2,\qquad
 \frac{\dot{\mathcal N}}{\mathcal N}=-\frac{\ii\hbar}{m_\parallel}a,
 \qquad
 \frac{\partial_\zeta^2\psi}{\psi}=-2a+4a^2\zeta^2.
\end{equation*}
TDSE 左右两侧的常数项与 $\zeta^2$ 项于是分别相等。再用
$\Re(a_0/d)=\Re a_0/|d|^2$，对 $|\psi|^2$ 作高斯积分，归一化保持为一。所以传播后的函数同时满足初始条件、运动方程和归一化条件。
\end{derivation}

% 数学结论及完整证明已移入正文相应公式之后。


% 数学结论及完整证明已移入正文相应公式之后。


% 数学结论及完整证明已移入正文相应公式之后。


% 数学结论及完整证明已移入正文相应公式之后。


% 数学结论及完整证明已移入正文相应公式之后。

